CH has been implicitly referenced in redistricting case law and has been
made explicit in several state redistricting standards.

\subsection{Shaw v.\ Reno (1993)}

The Supreme Court's visual-irregularity test in \textit{Shaw v.\ Reno},
509 U.S.\ 630 (1993), is a conceptual predecessor to CH~\citep{shaw1993}.
The Court held that a district whose shape ``is so bizarre on its face'' that
it cannot be explained except as a racial classification triggers strict scrutiny.
This test requires the court to examine whether the district's shape is the
result of a ``legitimate districting principle'' or of deliberate racial sorting.
CH provides a rigorous operationalisation: a district with CH $< 0.60$ has
more than 40\% of its convex hull unoccupied, which is the geometric signature
of a corridor or tentacle shape that cannot be explained by contiguity alone.

\subsection{Ohio Redistricting Criteria}

Ohio's 2015 redistricting reform (H.B.\ 1) requires that districts be
``as compact as possible'' using the principle that districts should not
``bypass populations'' on their way to connecting distant communities.
This ``bypass'' criterion is equivalent to a high-CH requirement:
a district that bypasses populations has a convex hull that is much larger
than the district itself (low CH).

\subsection{Pennsylvania: League of Women Voters v.\ Commonwealth}

In \textit{League of Women Voters of Pennsylvania v.\ Commonwealth},
178 A.3d 737 (Pa.\ 2018), the Pennsylvania Supreme Court struck down the
2011 congressional map as an unconstitutional partisan gerrymander under
the Pennsylvania Constitution.
The court's opinion cited the ``meandering and irregular'' shape of several
districts, which expert reports documented using CH-like measures~\citep{lwvpa2018}.
The court-ordered remedial map produced districts with substantially higher
CH values than the struck-down map.

\subsection{Recommended Disclosure}

For court filings using CH:
\begin{quote}
  \emph{All Convex Hull Ratio scores were computed using the convex hull of
  district polygons derived from 2020 Census TIGER/Line tract geometries
  projected to Albers Equal Area Conic (EPSG:5070). The convex hull was
  computed using Andrew's monotone chain algorithm. CH $= A(\text{district}) /
  A(\text{convex hull})$.}
\end{quote}
